\section{Scope and Related Work}

\subsection{Digital therapeutics and serious games for ADHD}

The STARS-ADHD trial and subsequent reviews support continued study of digital and game-based interventions in pediatric ADHD while also showing heterogeneity in intervention content, comparators, outcomes, and evidence for generalization \parencite{kollins2020stars,oh2024dtx,lin2025seriousgames,schultz2025content,tommasi2026serious}. ARC-ADHD therefore treats digital delivery as a means of delivering controlled practice, not as evidence of therapeutic efficacy by itself.

\subsection{Emotion regulation as a distinct target}

Emotion dysregulation is associated with ADHD and functional burden \parencite{graziano2016dysregulation}. Psychosocial interventions addressing dysregulation are heterogeneous, and executive-control training cannot be assumed to produce emotion-regulation change merely because both involve control processes \parencite{vacher2020psychosocial,groves2022emotion}. The present model therefore requires an independently assessed emotion-regulation or dysregulation outcome in addition to in-game performance.

\subsection{Why recovery is treated as a temporal target}

Emotion regulation is not only a static level of distress, strategy endorsement, or informant-rated difficulty. Developmental emotion research also characterizes within-person temporal dynamics, including intensity, variability, instability, inertia, and change over time \parencite{reitsema2022emotiondynamics}. In developmental psychobiology, structured self-regulation paradigms explicitly distinguish baseline, challenge reactivity, within-task change, and post-challenge recovery as different temporal contexts \parencite{davis2020context}. This provides a theoretical rationale for making return to goal-directed behavior after a controlled perturbation the central proximal quantity in ARC-ADHD.

This rationale does not establish \(T_{\mathrm{rec}}\) as a validated clinical measure. A recent systematic review of behavioral and physiological emotion-regulation assessment in children and adolescents found substantial heterogeneity in definitions, theoretical embedding, and psychometric reporting, including limited reporting of validity evidence \parencite{bunz2026ermeasurement}. ARC-ADHD therefore treats recovery latency as a project-specific mechanistic measure whose construction, stability, sensitivity to missingness, challenge dependence, and relation to independent outcomes must be established empirically before strong inference.

\subsection{Why the core paper is deliberately narrow}

Earlier versions of ARC-ADHD combined multisensory cueing, sensorimotor synchronization, body awareness, augmented reality, facial behavior, and bilateral audio in a single framework. Such modularity is useful for product design but problematic for causal inference because each component creates distinct mechanisms, measurement error, safety constraints, and statistical contrasts.

The revised programme retains only components required to instantiate the challenge-regulation-recovery cycle in the core paper. Visual, auditory, and haptic cues may be used as bounded information channels, and a brief body-orientation step may be used to connect strategy selection with subjective state. Stronger claims about movement, sensorimotor synchronization, spatial context, and augmented reality are reserved for the embodied/AR companion study. Camera-derived facial behavior and bilateral auditory stimulation are fully outside the core intervention claim. Wearable EEG, including the Muse S Athena companion study, is likewise an optional measurement layer and is not required to define, deliver, or validate the core recovery event.

\subsection{Replicability of a complex behavioral intervention}

The intervention should be described using the TIDieR framework so that rationale, materials, procedures, provider role, delivery mode, setting, dose, tailoring, modifications, and fidelity can be reconstructed \parencite{hoffmann2014tidier}. A later randomized-trial protocol should additionally follow SPIRIT 2025 \parencite{chan2025spirit}. This paper instead defines the mechanism and the development gates that must be completed before such a protocol is credible.
